The important progress in deep learning has led to a massive interest for these approaches in industry. However, when applying machine learning to critical tasks such has in the transportation or the medical domain, empirical and theoretical guarantees are required. Unfortunately, it has been shown that neural networks are weak to adversarial attacks: a carefully chosen small shift to the input, usually indistinguishable from noise, can change the class prediction \cite{szegedy2013}.  This sensitivity to adversarial attacks is mainly due to the Lipschitz constant of a neural network which can be arbitrarily high when unconstrained. Most of white-box attacks (where the full model is available) take advantage of it to build adversarial examples by using gradient descent with respect to the input variables. $FGSM$~\cite{Goodfellow2014a} performs only one step of gradient descent when other approaches such as $PGD$~\cite{Madry2017,carlini2017towards} find the optimal adversarial example iteratively. In black-box scenarios, gradients or logits of the model are not available. In such case, attacks start from large perturbations and then reducing it step by step (see for instance, boundary attacks~\cite{Brendel2018} and pointwise attacks~\cite{Schott2018}).

There are three major types of strategy to address the issue of adversarial attacks. Agnostic defenses are independent of the model and consist in altering the input or the prediction. For instance, Cohen et al. obtain a provable certificate with respect to $l_2$ norm by using Gaussian random smoothing~\cite{pmlr-v97-cohen19c}. DEFENSE-GAN~\cite{samangouei2018defensegan} uses a GAN to transform the input into the closest non-adversarial one at inference time. The second group of strategies relies on saddle point optimization by adding a penalty term measuring the empirical weakness against adversarial example during the learning process~\cite{Madry2017}. The last type of approaches focus on the Lipschitz constant of the network. It has been proven that bounding the  Lipschitz constant of a neural network provides certifiable robustness guarantees against local adversarial attacks~\cite{hein_formal_2017,ono_lightweight_2018}, improves generalizations~\cite{Sokolic_2017} and the interpretability of the model~\cite{tsipras2018robustness}.  This constraint can be achieved layer by layer by using spectral normalization~\cite{cisse_parseval_2017} or non-expansive layer~\cite{qian_l2-nonexpansive_2019}. In \cite{li2019preventing}, Li et al. go beyond the Lipschitz constant bounding by requiring layers to be gradient norm preserving. Combined with hinge loss, it allows them to achieve stronger robustness certification. However, the main limitation of this approach relies in the link between the hinge margin parameter and the robustness of the network.

In this paper we propose a new classification framework based on optimal transport that integrates the Lipschitz constant and the gradient norm preserving constraint as a theoretical requirement. To the best of our knowledge, very few researches investigate the link between binary classification and optimal transport (in~\cite{Frogner2015}, Frogner et al. use Wasserstein loss to improve multilabel classification). In Wasserstein GAN~\cite{Arjovsky2017}, the k-Lipschitz networks used to measure the distance between two distributions act like a discriminator, in analogy with the initial GAN algorithm \cite{Goodfellow2014}. The Wasserstein distance is approximated using a loss based on the Kantorovich-Rubinstein dual formulation and a k-Lipschitz network constrained by weight clipping. However, as we will demonstrate, a vanilla classifier based on the Kantorovich-Rubinstein loss is suboptimal, even on toy datasets.

We propose a binary classifier based on a regularized version of Kantorovich-Rubinstein formulation using a hinge loss term. We show that it remains the dual of an optimal transport problem, and we prove that the optimal solution of the problem exists and makes no error when the classes are well separated. With this new optimization problem, we guarantee to have an accurate classifier with a loss that is defined on and takes advantage of 1-Lipschitz function. As in~\cite{li2019preventing}, we bound the Lipschitz constant of the linear layers by B\"jorck normalization and use norm preserving activation functions~\cite{pmlr-v97-anil19a}. However, the optimal transport interpretation of the problem makes the bridge between these constraints and the loss function. When solving this optimal transport problem, attacking a prediction corresponds to travel along the transport plan up to the decision frontier. The output of the optimal network is linked to the length of this path, which is maximized during the optimization process of the proposed loss. 

The paper, and the contributions, are structured as follows. In Section~\ref{sec:wass_distance}, we recall the definition of Wasserstein distance and the dual optimization problem associated. We present the interesting properties of a classifier based on this approach, illustrate that it leads to a suboptimal classifier even on a toy dataset. Section~\ref{sec:kr_hinge} describes the proposed binary classifier, based on a regularized version of the Kantorovich-Rubinstein loss with a hinge regularization. We show that the primal of this classification problem is a new optimal transport problem and we demonstrate different mathematical properties of our approach. Section~\ref{sec:lip_const} is devoted to the way of constraining the classifier to be  1-Lipschitz and how to generalize the approach to multi classification problems. Section~\ref{sec:experimentation} presents the results of  experiments on MNIST, Cifar10 and CelebA datasets, measuring and comparing the results of different approaches in terms of accuracy and robustness. Last, we demonstrate that with our approach, building an adversarial example requires explicitly changing the example to an in-between two-classes image, which correspond to a point halfway on the transport plan. Proofs, computations details and additional experiments are reported in the appendix. 
